% GENERATED FILE. Edit canonical lessons, then run npm run generate:books.
% canonical-source-hash: fnv1a64:a05d529e9b3f72ed
% canonical-lessons: ZH-C74-jin, ZH-W74-jin-enter, ZH-C74-keren, ZH-W74-ke-guest, ZH-C74-menkou, ZH-C74-huanying, ZH-C74-jieshao, ZH-R74-review-70-74

\chapter{Visiting}
\label{ch:zh-visiting}
\hlchaptermodality{74}

\begin{chapteropening}
\textbf{By the end of this chapter:} \emph{I can ask a guest in, welcome them and introduce people.}

Complete the last lesson of chapter 74: I can ask a guest in, welcome them and introduce people.
\end{chapteropening}

\section[\texorpdfstring{\zh{进} (jìn)}{进 (jìn)}]{\zh{进} (jìn) — to come in; to enter}

\label{lesson:ZH-C74-jin}



\begin{quote}
Before the new one: say the Mandarin for airport, then the Mandarin for ticket.
\end{quote}

\subsection*{You'll want to know: \zh{进}}
\textbf{\zh{进}} — \emph{jìn} — \textquotedblleft{}to come in; to enter\textquotedblright{}.

\begin{quote}
\textbf{\zh{他进门了}} \emph{tā jìn mén le} — He's come in at the door.
\end{quote}

The new character, in pieces: \zh{进} — a four-stroke piece read jǐng, a well, on the road piece of \zh{这}. The next lesson puts it in your hand.

\begin{cousinweb}
\textbf{\zh{进}} (\emph{jìn}) is inherited from Middle Chinese \emph{tsinH} (Guangyun, in Baxter's transcription), glossed 'forward' and 'to advance'.
\end{cousinweb}

\begin{grammarlens}[title={What you've built}]
One more word for meeting people.
\end{grammarlens}

\subsection*{Your turn}
\begin{itemize}
\raggedright
  \item \emph{Say it:} \emph{jìn}
  \item \emph{Say it:} \emph{jìn}, once more
  \item \emph{Say it:} \emph{tā jìn mén le} — He's come in at the door.
  \item \emph{Recall:} say the Mandarin for airport, then the Mandarin for ticket, then say \emph{jìn} again
\end{itemize}

\begin{tcolorbox}[breakable,colback=teal!4,colframe=teal!35!black,title={Before you move on}]
What does \zh{进} mean? (\textquotedblleft{}To come in; to enter\textquotedblright{}.) Say it once more.
\end{tcolorbox}
\section[\texorpdfstring{\zh{进} (jìn)}{进 (jìn)}]{\zh{进} — jǐng on the road piece}

\label{lesson:ZH-W74-jin-enter}



\begin{quote}
The word first, then the character inside it.

\begin{itemize}
\raggedright
  \item \emph{Recall:} say \emph{jìn} — \textbf{R1}, one lesson back, before it decays
\end{itemize}

\begin{quote}
\zh{开} \zh{这}
\end{quote}

\emph{Read it:} each one aloud, with nothing said first — \emph{kāi}, \emph{zhè}
\end{quote}

\subsection*{Script — jǐng on the road piece}
\hlblockfigure{\detokenize{figures/ZH-W74-jin-enter-filmstrip.pdf}}{How \zh{进} is written, stroke by stroke}

\begin{quote}
\zh{进}
\end{quote}

\textbf{Two pieces: jǐng first, the road last}, though the road runs up the left and under.

\textbf{Jǐng, four strokes}: two horizontals, then a left-falling stroke and a vertical down through them, the same four as \zh{开}.

\textbf{The road piece, three strokes}, as in \zh{这}: a dot, then a stroke that turns and turns again, then a long sweep along the bottom to the right.

4 + 3 = seven.

\subsection*{Your turn}
\begin{itemize}
\raggedright
  \item \emph{Write it:} \zh{进}, all seven, with the model in front of you
  \item \emph{Cover:} write it again from memory, and count the pieces, not the strokes
\end{itemize}

\begin{tcolorbox}[breakable,colback=teal!4,colframe=teal!35!black,title={Before you move on}]
How many pieces in \zh{进}? (\textbf{Two}.) Which character you know has the same first four strokes? (\textbf{\zh{开}}.) Which character has the same road piece? (\textbf{\zh{这}}.)

Source: \href{https://raw.githubusercontent.com/chanind/hanzi-writer-data/68d10a4b21150cae5e1ebbd223eed289cf32d90c/data/\%E8\%BF\%9B.json}{Hanzi Writer Data \zh{进}, pinned stroke paths and medians}.
\end{tcolorbox}
\section[\texorpdfstring{\zh{客人} (kèrén)}{客人 (kèrén)}]{\zh{客人} (kèrén) — guest}

\label{lesson:ZH-C74-keren}



\begin{quote}
Before the new one: say the Mandarin for ticket, then the Mandarin for to come in.
\end{quote}

\subsection*{You'll want to know: \zh{客人}}
\textbf{\zh{客人}} — \emph{kèrén} — \textquotedblleft{}guest\textquotedblright{}.

\begin{quote}
\textbf{\zh{客人来了快开门}} \emph{kèrén lái le kuài kāi mén} — The guests are here: quick, open the door!
\end{quote}

The new character, in pieces: \zh{客} — the roof of \zh{家}, a walking foot under it, and \zh{口} at the bottom. The next lesson puts it in your hand.

\begin{cousinweb}
\textbf{\zh{客人}} (\emph{kèrén}) is built from \zh{客} 'guest' and \zh{人} 'person', both inherited from Middle Chinese, \emph{khaek} and \emph{nyin} (Guangyun, in Baxter's transcription).
\end{cousinweb}

\begin{grammarlens}[title={What you've built}]
One more word for meeting people.
\end{grammarlens}

\subsection*{Your turn}
\begin{itemize}
\raggedright
  \item \emph{Say it:} \emph{kèrén}
  \item \emph{Say it:} \emph{kèrén}, once more
  \item \emph{Say it:} \emph{kèrén lái le kuài kāi mén} — The guests are here: quick, open the door!
  \item \emph{Recall:} say the Mandarin for ticket, then the Mandarin for to come in, then say \emph{kèrén} again
\end{itemize}

\begin{tcolorbox}[breakable,colback=teal!4,colframe=teal!35!black,title={Before you move on}]
What does \zh{客人} mean? (\textquotedblleft{}Guest\textquotedblright{}.) Say it once more.
\end{tcolorbox}
\section[\texorpdfstring{\zh{客} (kè)}{客 (kè)}]{\zh{客} — under the roof, a foot and \zh{口}}

\label{lesson:ZH-W74-ke-guest}



\begin{quote}
The word first, then the character inside it.

\begin{itemize}
\raggedright
  \item \emph{Recall:} say \emph{kèrén} — \textbf{R1}, one lesson back, before it decays
\end{itemize}

\begin{quote}
\zh{家} \zh{字} \zh{口}
\end{quote}

\emph{Read it:} each one aloud, with nothing said first — \emph{jiā}, \emph{zì}, \emph{kǒu}
\end{quote}

\subsection*{Script — under the roof, a foot and \zh{口}}
\hlblockfigure{\detokenize{figures/ZH-W74-ke-guest-filmstrip.pdf}}{How \zh{客} is written, stroke by stroke}

\begin{quote}
\zh{客}
\end{quote}

\textbf{Three pieces, top to bottom.}

\textbf{The roof, three strokes}, as on \zh{家} and \zh{字}: a dot, a short stroke down the left, and a horizontal that turns down and hooks.

\textbf{The walking foot, three strokes}: a short left-falling stroke, a horizontal that turns into a long left-falling stroke, and a long stroke falling to the right.

\textbf{\zh{口}, three strokes}, at the bottom.

3 + 3 + 3 = nine.

\subsection*{Your turn}
\begin{itemize}
\raggedright
  \item \emph{Write it:} \zh{客}, all nine, with the model in front of you
  \item \emph{Cover:} write it again from memory, and count the pieces, not the strokes
\end{itemize}

\begin{tcolorbox}[breakable,colback=teal!4,colframe=teal!35!black,title={Before you move on}]
How many pieces in \zh{客}? (\textbf{Three}.) Which characters you know have the same roof? (\textbf{\zh{家} and \zh{字}}.) What sits at the bottom? (\textbf{\zh{口}}.)

Source: \href{https://raw.githubusercontent.com/chanind/hanzi-writer-data/68d10a4b21150cae5e1ebbd223eed289cf32d90c/data/\%E5\%AE\%A2.json}{Hanzi Writer Data \zh{客}, pinned stroke paths and medians}.
\end{tcolorbox}
\section[\texorpdfstring{\zh{门口} (ménkǒu)}{门口 (ménkǒu)}]{\zh{门口} (ménkǒu) — doorway; entrance}

\label{lesson:ZH-C74-menkou}



\begin{quote}
Before the new one: say the Mandarin for to come in, then the Mandarin for guest.
\end{quote}

\subsection*{You'll want to know: \zh{门口}}
\textbf{\zh{门口}} — \emph{ménkǒu} — \textquotedblleft{}doorway; entrance\textquotedblright{}.

\begin{quote}
\textbf{\zh{门口有一个客人}} \emph{ménkǒu yǒu yí ge kèrén} — There's a guest at the door.
\end{quote}

\begin{cousinweb}
\textbf{\zh{门口}} (\emph{ménkǒu}) is built from \zh{门} 'door' and \zh{口} 'mouth', both inherited from Middle Chinese, \emph{mwon} and \emph{khuwX} (Guangyun, in Baxter's transcription): the mouth of the door, where you go in.
\end{cousinweb}

\begin{grammarlens}[title={What you've built}]
One more word for meeting people.
\end{grammarlens}

\subsection*{Your turn}
\begin{itemize}
\raggedright
  \item \emph{Say it:} \emph{ménkǒu}
  \item \emph{Say it:} \emph{ménkǒu}, once more
  \item \emph{Say it:} \emph{ménkǒu yǒu yí ge kèrén} — There's a guest at the door.
  \item \emph{Recall:} say the Mandarin for to come in, then the Mandarin for guest, then say \emph{ménkǒu} again
\end{itemize}

\begin{tcolorbox}[breakable,colback=teal!4,colframe=teal!35!black,title={Before you move on}]
What does \zh{门口} mean? (\textquotedblleft{}Doorway; entrance\textquotedblright{}.) Say it once more.
\end{tcolorbox}
\section[huānyíng]{huānyíng — welcome}

\label{lesson:ZH-C74-huanying}



\begin{quote}
Before the new one: say the Mandarin for guest, then the Mandarin for doorway.
\end{quote}

\subsection*{You'll want to know: huānyíng}
\emph{huānyíng} — \textquotedblleft{}welcome\textquotedblright{}.

\begin{quote}
\emph{huānyíng nǐ lái wǒ jiā} — Welcome to my home!
\end{quote}

Its characters come later; for now it is a word for the ear.

\begin{cousinweb}
\emph{huānyíng} is built from huān 'glad' and yíng 'to go to meet', both inherited from Middle Chinese, \emph{xwan} and \emph{ngjaeng} (Guangyun, in Baxter's transcription): glad to meet you.
\end{cousinweb}

\begin{grammarlens}[title={What you've built}]
One more word for meeting people.
\end{grammarlens}

\subsection*{Your turn}
\begin{itemize}
\raggedright
  \item \emph{Say it:} \emph{huānyíng}
  \item \emph{Say it:} \emph{huānyíng}, once more
  \item \emph{Say it:} \emph{huānyíng nǐ lái wǒ jiā} — Welcome to my home!
  \item \emph{Recall:} say the Mandarin for guest, then the Mandarin for doorway, then say \emph{huānyíng} again
\end{itemize}

\begin{tcolorbox}[breakable,colback=teal!4,colframe=teal!35!black,title={Before you move on}]
What does huānyíng mean? (\textquotedblleft{}Welcome\textquotedblright{}.) Say it once more.
\end{tcolorbox}
\section[jièshào]{jièshào — to introduce}

\label{lesson:ZH-C74-jieshao}



\begin{quote}
Before the new one: say the Mandarin for doorway, then the Mandarin for welcome.
\end{quote}

\subsection*{You'll want to know: jièshào}
\emph{jièshào} — \textquotedblleft{}to introduce\textquotedblright{}.

\begin{quote}
\emph{qǐng nǐ jièshào nǐ de jiārén} — Tell us about your family, please.
\end{quote}

Its characters come later; for now it is a word for the ear.

\begin{cousinweb}
\emph{jièshào} is built from jiè and shào, both inherited from Middle Chinese, \emph{keajH} and \emph{dzyewX} (Guangyun, in Baxter's transcription). The dictionary glosses the first as 'to help' and the second as 'to carry on'; how the pair came to mean introducing somebody is not settled.
\end{cousinweb}

\begin{grammarlens}[title={What you've built}]
One more word for meeting people.
\end{grammarlens}

\subsection*{Your turn}
\begin{itemize}
\raggedright
  \item \emph{Say it:} \emph{jièshào}
  \item \emph{Say it:} \emph{jièshào}, once more
  \item \emph{Say it:} \emph{qǐng nǐ jièshào nǐ de jiārén} — Tell us about your family, please.
  \item \emph{Recall:} say the Mandarin for doorway, then the Mandarin for welcome, then say \emph{jièshào} again
\end{itemize}

\begin{tcolorbox}[breakable,colback=teal!4,colframe=teal!35!black,title={Before you move on}]
What does jièshào mean? (\textquotedblleft{}To introduce\textquotedblright{}.) Say it once more.
\end{tcolorbox}
\section[(dialogue)]{Looking back: from More to Eat to Visiting}

\label{lesson:ZH-R74-review-70-74}



\begin{quote}
Say the Mandarin for meat and steamed bun.
\end{quote}

\subsection*{Your turn}
\begin{itemize}
\raggedright
  \item \emph{Say it:} meat, steamed bun, noodles, dumpling, soup
  \item \emph{Say it:} nice to drink, sweet, spicy, salty, to order food
  \item \emph{Say it:} train, to drive, aeroplane, road, to arrive
  \item \emph{Say it:} to get off, to get on, railway station, airport, ticket
  \item \emph{Say it:} to come in, guest, doorway, welcome, to introduce
  \item \emph{Read it:} \zh{火} \zh{开} \zh{下} \zh{进} \zh{客} — each one aloud, then a word that holds it
\end{itemize}

\begin{tcolorbox}[breakable,colback=teal!4,colframe=teal!35!black,title={Before you move on}]
Meat? (\emph{ròu}.) Steamed bun? (\emph{bāozi}.) Noodles? (\emph{miàntiáo}.) Dumpling? (\emph{jiǎozi}.) Soup? (\emph{tāng}.) Nice to drink? (\textbf{\zh{好喝}}, \emph{hǎohē}.) Sweet? (\emph{tián}.) Spicy? (\emph{là}.) Salty? (\emph{xián}.) To order food? (\emph{diǎncài}.) Train? (\textbf{\zh{火车}}, \emph{huǒchē}.) To drive? (\textbf{\zh{开车}}, \emph{kāichē}.) Aeroplane? (\emph{fēijī}.) Road? (\emph{lù}.) To arrive? (\emph{dào}.) To get off? (\textbf{\zh{下车}}, \emph{xiàchē}.) To get on? (\textbf{\zh{上车}}, \emph{shàngchē}.) Railway station? (\emph{huǒchēzhàn}.) Airport? (\emph{jīchǎng}.) Ticket? (\emph{piào}.) To come in? (\textbf{\zh{进}}, \emph{jìn}.) Guest? (\textbf{\zh{客人}}, \emph{kèrén}.) Doorway? (\textbf{\zh{门口}}, \emph{ménkǒu}.) Welcome? (\emph{huānyíng}.) To introduce? (\emph{jièshào}.)
\end{tcolorbox}
